% !TeX root = ../main.tex

Das Nutzen einer Datenbank bringt viele Vorteile, aber auch Nachteile.
Vor allem, wenn die Fehlerbehandlung nicht solide läuft.
Denn wenn eine Datenbank einen Fehler bemerkt, wird dieser bisher manuell behandelt.
Sobald ein Fehler auftritt, bricht die Datenbank die Transaktion ab und gibt dem Client den Fehlercode zurück.
In dem Fall muss dann die Clientapplikation entscheiden wie mit diesem um gegangen wird.
Das heißt meistens, dass der User eigenhändig eingreifen muss.
Das Ziel ist es, dass der Client schon gar nicht mehr erfährt, dass eine Transaktion abgebrochen wurde.
Die folgende Arbeit behandelt, wie genau man bestimmte Fehler abfangen und möglicherweise auch lösen kann.
\textbf{Muss auch dringen noch geändert werden.}